\documentclass[11pt,a4paper]{article}
\usepackage[utf8]{inputenc}
\usepackage[T1]{fontenc}
\usepackage[margin=2.3cm]{geometry}
\usepackage{booktabs,longtable,array,amsmath,enumitem}
\usepackage{xurl}
\usepackage[hidelinks]{hyperref}
\setlist{itemsep=1pt,topsep=3pt}
\newcolumntype{L}[1]{>{\raggedright\arraybackslash}p{#1}}
\newcommand{\est}{\textsc{est}}
\newcommand{\ver}{\textsc{ver}}
\title{HL-LHC electroweak quirks, study v2\\[2pt]\large Phase 1 report: strategy, feasibility and Gate A}
\author{Private work (simulation). Public information only.}
\date{16 November 2026 (day 0 of the approved calendar)}
\begin{document}
\maketitle

\noindent\textbf{Tags.} \ver{} = verified in this study (run or read, with file and line in the repository).
\est{} = planning estimate. Every reach and background number in this report is an \est{} built on the placeholder
efficiencies of v1, and no v2 signal efficiency has been computed.

\section{Summary}
\begin{itemize}
\item Phase 1 is finished. The strategy passed its review panel in the tenth cycle, and Gate A was passed on
      16 November 2026 with every default answer.
\item Approved schedule: 4{,}609 CPU-h, stop-and-ask at 5{,}992 CPU-h, about 60 days. The original target of
      2--4 weeks is missed by 4--5 weeks.
\item The v1 chain runs on the Linux cluster, the event generators and Geant4 are installed, and a resumable job
      runner is validated (\ver).
\item Expected outcome of v2 (\est): a background that is estimated and tested, of about
      $0.7$--$1.0\times10^{4}$ events led by photon conversions from the hard scatter, and a reach below 243~GeV.
      Existing bounds are at about 180~GeV. The upside, 710--790~GeV, is not the expectation.
\item The background projection of day 15 (1 December) is expected to exceed its threshold. The study then
      pauses and asks how to continue.
\end{itemize}

\section{What v1 was, and why v2 exists}
The search looks for a quirk--antiquirk pair bound by a string. In a tracker its hits do not form helices. They lie
in a thin plane through the collision vertex. A dedicated plane finder searches for this pattern after standard
tracking, in events triggered by missing transverse momentum, at 14~TeV, 3~ab$^{-1}$ and pile-up 200. Selections
form a ladder of 23 tightening cuts, R01--R23.

v1 made no physics statement:
\begin{itemize}
\item No selection met the frozen background rule. The lowest background figure was $1.23\times10^{6}$ events,
      about seven decades above the budget of 0.11.
\item From R08 on, every leading background term was a zero-count bound: no simulated event passed, so the term
      measured the size of the simulated sample, not the background.
\item The pile-up extrapolation failed its own closure test (factors 1.50 and 1.61) and was patched by a factor
      2.41.
\item The best expected limit was 790 times the predicted cross section, so nothing was excluded.
\end{itemize}
v1 therefore showed that counting simulated events cannot prove that the background is small. It did not show
that the background is large. v2 replaces the bound by an estimate with closure tests and a stated width.

\section{Improvements of v2 with respect to v1}
The table lists every item of the publication plan as the approved strategy commits to it. Priority M is a
must-have item and S a should-have item. The acceptance criteria are fixed in advance in
\texttt{COMMITMENTS.md}; none has been evaluated yet.

\begin{longtable}{L{1.0cm}L{0.5cm}L{5.4cm}L{7.6cm}}
\toprule
Item & Pr. & In v1 & In v2 \\
\midrule
\endhead
\bottomrule
\endfoot
C1 & M & Ran on one machine; some files lost; no independent rerun.
   & v1 is rerun on the cluster from logged port edits: 41 declared comparisons with v1, each with a
     pre-declared test. Tracked lockfiles and a container recipe. Six of ten signal counts stay sealed until Gate B. \\
C2 & M & Limited validation of the plane finder.
   & Brute-force cross-check with new code on 50 signal and 50 pile-up trials (agreement within 1~mrad in at least
     48 of 50), bit-identical reruns, and a seed-threshold stability test. \\
C3 & M & Internal labels, unsourced statements and errors in the paper text.
   & All 24 listed corrections, an errata file for v1, novelty wording checked against primary sources, and
     disclosure of the AI workflow. \\
B1 & M & Decays in flight and nuclear interactions from generator-level rates; the decay term rests on 1--2 counts.
   & Geant4 vertices handed to the fast tracker simulation. An estimator built from pass and near-pass
     probabilities without pile-up, a survival factor and a promotion model, with held-out closure within a
     factor 2. A no-credit ceiling for every factor that cannot be tested, itself tested. Geant4 geometry matched
     to the tracking geometry within 1\,\% of a radiation length per layer. \\
B2 & M & Thinning extrapolation of pile-up failed closure and was patched by $\times2.41$.
   & Repaired thinning with no scale-factor patch, plus an independent semi-analytic occupancy estimator in new
     code. Each must reproduce direct counts within a factor 1.5 on held-out rungs, and the two must agree beyond
     them. A direct sample of 24{,}000 events gives about 70 expected counts at R07. \\
B3 & M & No tested conditional probabilities for the purpose-built cuts.
   & Each conditional is measured after at least three preceding rungs and tested on disjoint halves. A
     displaced-vertex veto is designed on background only. A time tag earns credit only where it closes and only
     with sourced timing tails. A failed or untested conditional earns no credit. \\
B4 & M & Eight Standard Model processes entered by proxy; Z($\tau\tau$), lost-lepton Z($\ell\ell$) and multijets
         were absent.
   & Hit-level samples for the proxied processes, the three missing processes added with a missing-momentum model
     built from cited public parametrisations and validated against a published figure, cited K-factors, and a
     generator comparison as normalisation uncertainty. \\
B5 & S & No data-driven estimate.
   & A road-rotation sideband closed in simulation, with its flatness, non-flat part and road independence
     measured. Reduced form approved: 8 roads on 6{,}000 events. \\
B6 & S & One detector model, no variations.
   & Seven detector-model variations as robustness tests (resolution, inefficiency, material, noise top-up,
     generator tune). Nuisances enter the fit only with a cited size. Reduced form approved: 4 signal points. \\
B7 & S & Cosmic muons and beam backgrounds not addressed.
   & A number from public rates, pointing and timing; a term of the fit if above 10\,\% of the background
     uncertainty. \\
T1 & M & One hybrid-toy limit; zero counts treated ad hoc.
   & A published \texttt{pyhf} workspace with one auxiliary Poisson channel per simulated count, a zero-count
     convention with a coverage test, closure-derived systematics, signal injection, pulls, impact ranking, and
     toy-based limits everywhere. It must reproduce v1's limit of 789.7 within 10\,\% in v1 mode. \\
T2 & M & A frozen background-free rule that no selection met, then a fallback selection.
   & Sixteen candidate selections frozen at Gate A, a figure of merit evaluated on an optimisation half, and
     results quoted on a disjoint evaluation half split by a salted hash. Bias checks on both the signal and the
     background side. An always-eligible baseline that uses direct counts only. \\
S1 & M & 200 signal rows per point; a fast path at some points.
   & 2{,}000 rows at each of the 40 points through the full chain, no fast path, efficiencies blind until the
     closure verdicts. \\
S2 & M & Leading-order cross sections; two definitions of the generator-cut fraction.
   & NLO K-factors after three validations ($K=1.165\pm0.002$ at 200~GeV in Phase 1, \ver), scale and PDF
     errors with two PDF sets, a resummed normalisation as cross-check, one definition of the generator-cut
     fraction with its uncertainty, QED final-state radiation in the nominal shower, a trigger turn-on from a
     public curve, and the infracolour Coulomb correction. \\
P1 & M & Thickness law with an untested numerical floor and one failed gate.
   & Extended-precision and mantissa-truncation floor tests, a paired field-off test with stated power, a
     pre-defined population estimator on new events, and the validity window stated as measured. \\
S3 & -- & -- & Endcap dropped by decision; barrel only. The endcap gain is stated as an outlook estimate: signal
     $\times1.5$--$3.2$ (\est). \\
S4, S5 & -- & -- & Not included (approved at Gate A). \\
\end{longtable}

\noindent Improvements that are not plan items:
\begin{itemize}
\item \textbf{Pre-registration.} Selections, criteria, estimators and tolerances are written, reviewed, hashed and
      committed before any production sample (planned for days 17--19).
\item \textbf{Blinding.} No signal efficiency is looked at before the background closure verdicts are recorded.
\item \textbf{Verdict classes.} Every criterion ends as pass, fail, not tested or not evaluable, with a fixed
      consequence. A failed test is reported and not tuned away.
\item \textbf{Reported target.} The width of the background estimate, the share of zero-count terms and a label
      per term (no credit, bound, low count, estimate) are part of the result.
\item \textbf{Method cross-check.} A boosted-decision-tree selection with a pre-registered recipe is compared
      with the cut ladder at equal background.
\item \textbf{Comparison with the literature.} Four reference analyses are tabulated and their published numbers
      overlaid on the limit and reach figures. Twenty retrieval tasks take the external numbers from public sources, each logged; nothing is
      quoted from memory.
\item \textbf{Reproduction package.} Code, configuration, efficiency tables and the likelihood, together with a
      full analysis note.
\item \textbf{Referee questions.} Twenty-five anticipated questions are answered with a pointer to the text.
\end{itemize}

\section{What Phase 1 established (\ver)}
\begin{center}
\begin{tabular}{L{7.5cm}L{6.5cm}}
\toprule
Quantity & Measured on the cluster \\
\midrule
Full chain, one pile-up-200 event (4 jobs per node) & 13.5 s \\
Plane finder, one trial & 11.7 s \\
Pile-up density $\rho=2$, one event & 206 s, 1.46 GiB \\
String propagation per event at 1 / 3 / 5 / 10 keV & 0.62 / 18.5 / 67 / 35 s \\
Leading-order generation chain & about 34 ms per event \\
Likelihood fit; 1{,}000 toys & 0.12 s; 652 s \\
Job overhead, workflow engine against direct launch & 17 s against 2.4 s \\
NLO K-factor at 200 GeV & $1.165\pm0.002$ \\
\bottomrule
\end{tabular}
\end{center}
\begin{itemize}
\item MG5\_aMC 2.9.22 with Pythia 8.313 reproduces v1's leading-order cross section. Geant4 11.4.3 is installed
      on each worker node.
\item The no-pile-up event cost (45 ms) is a guess. Of the 4{,}453 CPU-h of work, 753 rest on measured unit
      costs, 2{,}322 on derived ones and 1{,}378 on guesses. The guesses are measured in the first week of Phase 2.
\end{itemize}

\section{Decisions fixed at Gate A}
\begin{center}
\begin{tabular}{L{1.0cm}L{13.2cm}}
\toprule
Q1 & Of the ten v1 signal counts used in the reproduction test, six stay sealed until Gate B and four are open. \\
Q2 & If the day-15 projection exceeds 300 background events, the study pauses and asks. \\
Q3 & A must-have closure criterion that cannot be tested enters at its ceiling, with no credit. \\
Q4 & 4{,}609 CPU-h and about 60 days; stop-and-ask at 5{,}992 CPU-h; reduced forms of B5 and B6; S4 and S5 out;
     the second hadronisation model at generator level only. \\
Q5 & Truth-level signal generation may run before the pre-registration hash. \\
\bottomrule
\end{tabular}
\end{center}
Eleven further defaults stand without change. The one with the most weight fixes how the day-15 rule reads
Standard Model processes that have no sample by then: all of them are read through the single Z($\nu\nu$) pilot
count, which puts the limit at about 12 counts (\est).

\section{Expected outcome (\est)}
\begin{itemize}
\item \textbf{Expected.} Background of about $0.7$--$1.0\times10^{4}$ events at the R08--R09 selections. With
      the strategy's own normalisation width (0.17) the reach is below 243~GeV.
\item \textbf{Day 15.} The rule reads the pilot before track finding and is expected at
      $3\times10^{4}$--$3\times10^{5}$ events against the threshold of 300.
\item \textbf{Upside.} 710--790~GeV if the ceilings turn out empty. No route to it has been demonstrated. At v1's
      central size for decays in flight, no conversion rejection alone is enough.
\item \textbf{Weak inputs.} Decays in flight rest on 1--2 counts in v1, and the conversion estimate on an assumed
      pair efficiency. Both can move the result by a large factor in either direction.
\end{itemize}
If the expected case holds, the paper is the first measured background for this search, the reach at that
background, and a sideband recipe as the way forward. It is a methods and feasibility paper, clearly more solid
than v1, with a modest headline.

\section{Schedule}
\begin{center}
\begin{tabular}{L{3.4cm}L{2.6cm}L{8.2cm}}
\toprule
Date & Day & Milestone \\
\midrule
21 November & 5 & Reproduction of v1 and finder checks; CPU re-projection with measured costs \\
23 November & 7 & Preliminary background projection \\
1 December & 15 & Binding projection; pause expected \\
3--5 December & 17--19 & Pre-registration reviewed, hashed and frozen \\
about 4--6 January & 49--51 & Gate B, before signal efficiencies are looked at \\
about 13--15 January & 58--60 & Gate C, final paper \\
\bottomrule
\end{tabular}
\end{center}

\section{Review record}
The strategy went through ten review cycles (physics, critical and constructive reviewers, then an arbiter; the
tenth cycle was a narrow mechanical check by one reviewer and the arbiter).
\begin{center}
\begin{tabular}{lcccccccccc}
\toprule
Cycle & 1 & 2 & 3 & 4 & 5 & 6 & 7 & 8 & 9 & 10 \\
\midrule
Must resolve & 11 & 4 & 1 & 0 & 0 & 0 & 0 & 0 & 0 & 0 \\
Must fix & 24 & 14 & 15 & 5 & 3 & 1 & 4 & 1 & 1 & 0 \\
\bottomrule
\end{tabular}
\end{center}
Cycles 1--3 rebuilt the method. Cycles 4--5 made the summary honest about the expected background. Cycles 6--9
changed no design choice; they corrected how the day-15 decision is described, whose root cause was an ambiguous
rule that cycle 8 pinned down.

\section{Incidents}
\begin{itemize}
\item 13 November: stale file handles after a file rename on the shared data volume. Fixed by the administrator;
      directories on that volume are never renamed.
\item 13 November: the head node ran out of memory and all running agents were killed. State was reconciled; no
      data was lost. The orchestrator now runs under a 10~GB cap and nothing heavy runs on the head node.
\item 14--16 November: the shared data volume and remote login to the workers hung. The cause was not established
      by this study and no repair was attempted. The cluster was confirmed healthy on 16 November.
\item Times written in the log before 14 November 01:20 were estimates and drifted by up to seven hours. Commit
      times are authoritative.
\end{itemize}

\section{Now running (Phase 2, days 0--5)}
\begin{enumerate}
\item Reproduction of v1 on the cluster and the finder cross-checks, with timing of the full chain.
\item Theory inputs, thickness-law clean-ups, literature retrievals and the Herwig installation attempt.
\item The pilot samples that feed the day-15 projection, the likelihood builder, and the Geant4 geometry and
      hand-over.
\end{enumerate}

\section*{Files}
\begin{itemize}
\item Strategy: \path{analyses/quirks_v2/phase1_strategy/outputs/STRATEGY_oscar_2026-11-16_10-41.md}
\item Commitments and Gate A answers: \path{analyses/quirks_v2/COMMITMENTS.md}
\item Gate A package: \path{analyses/quirks_v2/gates/GATE_A.md}
\item Final verdict: \path{analyses/quirks_v2/phase1_strategy/review/arbiter/STRATEGY_ARBITER_v10_quentin_2026-11-16_10-51.md}
\item Log: \path{analyses/quirks_v2/LOG.md}
\end{itemize}
\end{document}
